\documentclass[11pt]{article}
\usepackage[T1]{fontenc}
\usepackage[utf8]{inputenc}
\usepackage{lmodern,microtype}
\usepackage[margin=1in]{geometry}
\usepackage{amsmath,amssymb,amsthm,mathtools}
\usepackage{booktabs,longtable,array,enumitem}
\usepackage{xcolor,hyperref,listings}
\hypersetup{colorlinks=true,linkcolor=blue!45!black,urlcolor=blue!45!black,citecolor=blue!45!black,pdftitle={Uniform introduction on raw terms: a smaller presentation of iPRA},pdfauthor={Research prepared for Vladimir Reshetnikov}}
\setlist{itemsep=3pt,topsep=5pt}
\lstset{basicstyle=\small\ttfamily,breaklines=true,columns=fullflexible,keepspaces=true,frame=single,rulecolor=\color{black!20}}
\newtheorem{theorem}{Theorem}[section]
\newtheorem{lemma}[theorem]{Lemma}
\newtheorem{proposition}[theorem]{Proposition}
\newtheorem{corollary}[theorem]{Corollary}
\theoremstyle{definition}\newtheorem{definition}[theorem]{Definition}
\theoremstyle{remark}\newtheorem{remark}[theorem]{Remark}
\newcommand{\G}{\mathcal G}
\newcommand{\Gn}{\mathcal G_{\mathrm n}}
\newcommand{\Gd}{\mathcal G_{\mathrm d}}
\newcommand{\EE}{\mathcal E}
\newcommand{\ipra}{\mathrm{iPRA}}
\newcommand{\code}[1]{\texttt{#1}}
\newcommand{\FV}{\operatorname{FV}}
\newcommand{\wt}{\omega}
\newcommand{\one}{S0}
\newcommand{\Cl}[2]{\langle #1;\,y.#2\rangle}
\title{Uniform introduction on raw terms\large\\A smaller presentation of intuitionistic\\primitive recursive arithmetic}
\author{Research prepared for Vladimir Reshetnikov}
\date{November 12, 2026}
\begin{document}
\maketitle
\begin{abstract}
We simplify the latest raw guarded presentation in ConciseIPRA by deleting
its exceptional rule for recognizing a naming formula. Introduction always
uses the newest numerical parameter, and uses zero when the ambient
numerical scope is empty. There is no supplied witness, witness-mode bit,
search, naming-pattern test, or computation normalization. We prove equality
of all open guarded consequence relations with the unrestricted-witness
calculus. Consequently the ordinary realization remains equiconsistent and
bi-interpretable with full first-order intuitionistic primitive recursive
arithmetic; both numerical domains are whole, equality is actual equality,
and both numerical comparison graphs are the identity.

The proof closes a bootstrap dependency. It first embeds the purely
equational iterator calculus using regenerated proof families, then
constructs predecessor without numerical coverage, and only afterwards
derives arbitrary names. The inherited scalar predecessor construction is
retained; an equational interpolation lemma removes its apparent reliance
on Boolean case distinctions. A tightly matched, complete Python checker
shrinks from 3,688 to 3,439 normalized bytes, from 1,159 to 1,073 significant
tokens, and from 1,339 to 1,248 AST nodes. This is a proved local
simplification, not a claim of an absolute smallest axiomatization or a
size record over all certificate formats in the repository.
\end{abstract}

\section{What is new, and what is being compared}
\label{sec:comparison}

The comparison snapshot is ProveIt commit
\begin{center}\small\ttfamily
9f9545f5ffd9e0392d4a36639503aca12383450a.
\end{center}
The inspected directory is \code{ConciseIPRA}~\cite{repo}.
Its main README describes 79 filed research packages. In addition, the
current \code{docs/incoming} directory contains five archives. We retrieved
and inspected those archives as well. This matters: two already eliminate
supplied introduction witnesses. The result below is compared with that
newer work, rather than with an obsolete witness-taking kernel.

The incoming \code{parameter\_intro} package already uses the newest
parameter, or zero, with full computation normalization and a
counter-aware successor-seeded recursor~\cite{parameter}. We do not claim
that a fixed newest witness is a new idea by itself. The incoming
\code{deterministic\_introduction} package works on raw scalar iterator
terms~\cite{deterministic}. Its introduction chooses a witness by the rule
\[
 w(B,\Pi)=
 \begin{cases}
 t,&B\text{ is literally }C(y,0,x)=C(y,0,t),\ t\text{ outer-scoped},\\
 \text{newest parameter in }\Pi,&\Pi\ne\varnothing,\\
 0,&\Pi=\varnothing.
 \end{cases}
\]
Our change deletes the first line and its entire recognizer. The remaining
rule is sufficient on the same raw terms and the same guarded formulas.
Every other instruction in the selected checker is unchanged.

This distinction is the concrete simplification. It removes a mathematical
exception and its implementation, while retaining the existing arithmetic
basis and avoiding a normalization layer. It also establishes that the
exception was not needed to bootstrap arbitrary names.

\begin{table}[h]
\centering
\begin{tabular}{lrrr}
\toprule
Complete acceptance module & Bytes & Tokens & AST nodes\\
\midrule
Incoming raw deterministic introduction & 3688 & 1159 & 1339\\
Uniform introduction, this package & 3439 & 1073 & 1248\\
Reduction & 249 & 86 & 91\\
\bottomrule
\end{tabular}
\caption{Matched complete programs, including input handling, syntax,
scope checks, all primitive rules, and the accept/reject wrapper.}
\label{tab:main-size}
\end{table}

The repository also has smaller programs using other interfaces. For
example, the incoming raw evaluation-discharge checker reports 1,892
bytes, its zero-normal companion 1,858 bytes, and the full-normal
parameter-introduction checker 1,925 bytes~\cite{evaluation,parameter}.
Compact public-goal codecs are smaller still. Thus Table~\ref{tab:main-size}
does not establish a corpus-wide size minimum. The result is a strict
reduction of the current raw deterministic axiomatization, with a matched
improvement in all three stated source-size coordinates.

\section{The complete axiomatization}
\label{sec:calculus}

\subsection{Terms, formulas, and scopes}

There is one numerical sort. The finite raw terms are
\[
 t ::= x\mid 0\mid St\mid I(n,a;u.h).
\]
Only $u$ is bound in $h$. The counter $n$, initializer $a$, and all other
parameters of $h$ are outside this binder. Substitution is capture avoiding.
No computation equation is built into term equality or substitution.
Define the abbreviation
\[
 C(v,a,b):=I(v,a;u.b),\qquad 1:=S0,
\]
where the displayed state variable is absent from $b$.

A formula is an equation $s=t$ or a paired formula $D_{x,y}B$. There is
one syntactic restriction: if $B$ is an equation immediately under $D$,
it must literally have the form
\[
 I(y,0;u.h)=I(y,0;v.k).                 \tag{Guard}
\]
The updates $h,k$ may depend on their state variables, on $x,y$, and on
outer parameters. If $B$ is compound, its root is another $D$, and its
own innermost equation must satisfy the same restriction. An equation at
the root of a formula is unrestricted. The designated contradiction is
$0=S0$.

A numerical scope is an ordered list
$\Pi=(a_1,\ldots,a_n)$ of distinct parameters. Define
\[
 \boxed{\ \wt(\Pi)=a_n\text{ if }n>0,\qquad
               \wt(\varnothing)=0.\ }                 \tag{1}
\]
The ordering is proof syntax; it imposes no numerical ordering on elements
of a model. Every rule specifies which parameters are added and when.

\subsection{Contexts and ordinary assumptions}

A context is a finite reusable list of clauses $\langle A;y.B(y)\rangle$,
with $y$ absent from $A$. Both components are canonical formulas in the
appropriate scopes. Its ordinary reading is
$A\to\forall y B(y)$. Clause indices are syntactic addresses, not
object-level numbers.

Set $H=(I(0,0;u.0)=0)$. An ordinary assumption $F$ abbreviates the
constant clause $\langle H;v.F\rangle$, with $v$ absent from $F$.
The formula $H$ is an instance of the first computation schema, not an
extra axiom. Consequently an ordinary assumption is retrieved by Use
and one computation proof. A closed proof begins in empty numerical
scope and with an empty clause context.

Write $\Pi;\Gamma\vdash_{\Gn} A$ for derivability. In displays where
the numerical scope is clear, we suppress $\Pi$.

\subsection{Two arithmetic schemas and four inference rules}

The complete nonlogical computation schemas, in every context, are
\[
 I(0,a;u.h)=a,                                         \tag{I0}
\]
\[
 I(Sn,a;u.h)=h[I(n,a;u.h)/u].                          \tag{IS}
\]
They range over all well-scoped finite terms of the grammar above.
The remaining primitive rules are exactly the following.

\paragraph{Use.}
For $\langle A;y.B\rangle\in\Gamma$ and any well-scoped term $t$,
\[
 \frac{\Gamma\vdash A}{\Gamma\vdash B(t)}.             \tag{Use}
\]

\paragraph{Transport.}
For a canonical one-parameter formula template $F(v)$,
\[
 \frac{\Gamma\vdash s=t\qquad\Gamma\vdash F(s)}
      {\Gamma\vdash F(t)}.                            \tag{Tr}
\]
The placeholder is a numerical variable. No substitution of proofs or
test of semantic equality occurs in this rule.

\paragraph{Open.}
\[
 \frac{\Pi;\Gamma\vdash D_{x,y}B(x,y)\qquad
       \Pi,b;\Gamma,\langle B(b,0);y.B(b,y)\rangle\vdash A}
      {\Pi;\Gamma\vdash A}.                          \tag{Open}
\]
Here $b$ is fresh for $\Pi,\Gamma,A$ and the fixed parameters of $B$.
The branch is checked in the scope with $b$ appended; the conclusion is
checked in the old scope, so $b$ cannot escape, including inside an unused
iterator body.

\paragraph{Uniform introduction.}
Put $w=\wt(\Pi)$ \emph{before} adding the step parameter. Let
\[
 \eta_B(z)=\begin{cases}z,&B\text{ is an equation},\\
                         0,&B\text{ is compound}.
             \end{cases}
\]
Then
\[
 \boxed{\quad
 \frac{\Pi,z;\Gamma,B(w,\eta_B(z))\vdash B(w,Sz)}
      {\Pi;\Gamma\vdash D_{x,y}B(x,y)}.
 \quad}                                               \tag{Intro}
\]
The parameter $z$ is fresh for the old scope, context, and fixed matrix
parameters. At an equation this is a progress step; at a compound
matrix it is a zero/successor case step. The matrix must satisfy Guard
when it is an equation. Introduction receives no witness data.

\begin{remark}[Exhaustiveness]
There are no further primitive logical or nonlogical axioms, successor
axioms, equality axioms, induction schemas, or inference rules. In
particular, reflexivity, ordinary hypothesis use, structural rules, and
the usual intuitionistic connectives are derived. The checker has six
instruction tags: the two computations and the four rules above.
\end{remark}

Let $\G$ denote exactly this calculus with an arbitrary well-scoped
witness $t$ in Intro. Let $\Gd$ denote the incoming calculus with its
literal-name exception. The inclusions
\[
 \Gn\ \subseteq\ \G
 \qquad\text{and}\qquad \Gd\ \subseteq\ \G
\]
are immediate on derivations. The substantial direction is completeness
of $\Gn$.

\section{The equivalence theorem and its proof discipline}
\label{sec:main-theorem}

\begin{theorem}[Exact open-consequence equivalence]
\label{thm:main}
For every ordered numerical scope $\Pi$, every well-scoped canonical
clause context $\Gamma$, and every canonical formula $A$,
\[
 \Pi;\Gamma\vdash_{\Gn}A
 \quad\Longleftrightarrow\quad
 \Pi;\Gamma\vdash_{\G}A
 \quad\Longleftrightarrow\quad
 \Pi;\Gamma\vdash_{\Gd}A.
\]
The translations are effective transformations of finite derivations;
the terms, formulas, and external assumptions are unchanged.
\end{theorem}

The easy directions use the respective witnesses as unrestricted
witnesses. The difficult direction is proved in
Sections~\ref{sec:hereditary}--\ref{sec:completeness}. It has the following
dependency order:
\begin{enumerate}
\item Establish replacement by uniformly regenerated proof families.
\item Embed the purely equational iterator fragment, before any names.
\item Obtain predecessor in that fragment, without numerical coverage.
\item Derive arbitrary names using the newest parameter.
\item Regenerate source derivations under arbitrary receiving scopes,
      recovering arbitrary-witness introduction.
\end{enumerate}
There is a real issue here. Adding a numerical parameter changes
$\wt(\Pi)$. Replaying a proof unchanged in a larger scope is generally
invalid. The argument never assumes that replay is a structural rule.
Likewise, it never substitutes an arbitrary term into a completed
restricted target proof. Where a source step needs a new instance, its
proof is regenerated in that instance.

\section{Hereditary proofs and the arithmetic bootstrap}
\label{sec:hereditary}

\begin{definition}[Hereditary derivability]
Write $\Gamma\Vdash A$ if, for every injective renaming of its displayed
parameters into a finite receiving numerical scope and every receiving
clause context containing the renamed $\Gamma$, there is an effectively
generated $\Gn$ proof of the renamed $A$. The generation may use the
receiving scope and context as inputs.
\end{definition}

This is a metatheoretic property of a family of finite proofs. It is not
an additional rule or oracle in the checker. Identity renaming and the
unchanged context show that hereditary derivability implies ordinary
derivability. Context membership and all the algebraic rules have
hereditary proofs.

\begin{lemma}[Replacement with hereditary donors]
\label{lem:hcut}
Suppose $\Gamma\Vdash A$ and $\Pi;\Delta,A\vdash_{\Gn}B$, where
$\Delta$ contains $\Gamma$ in the displayed scope and $A$ is added as
an ordinary assumption. Then $\Pi;\Delta\vdash_{\Gn}B$. The same
statement holds under renamed embeddings of $\Gamma$. Replacement is
valid below numerical binders; all donor hypotheses must remain available.
\end{lemma}
\begin{proof}
Induct on the receiving derivation. Preserve the exact numerical scope
at each of its nodes. At a Use of the constant clause for the removed
assumption, generate the donor proof in the numerical scope and context
of that leaf. All other Use addresses are adjusted for the deleted clause.
All other rule nodes retain their own scopes and parameters; in
particular their selected newest witness is unchanged. Below Open and
Intro the donor is generated in the extended scope, rather than obtained
by weakening an already generated proof. The recursion is on a finite
receiving proof and proves the assertion without a target numerical
weakening theorem.
\end{proof}

\subsection{Padding and cases without names}

For an equation $E=(s=t)$ put
\[
 G_yE=(C(y,0,s)=C(y,0,t)),\qquad
 L(E)=D_{a,v}G_vE,
\]
where $a,v$ are fresh for $E$. Both are abbreviations in the raw syntax.
The coordinate $a$ is unused. Computation and transport give
$G_0E$ and $G_{Sz}E\dashv\vdash E$. Therefore
\[
 L(E)\dashv\vdash E                                      \tag{2}
\]
uniformly in every receiving scope and additional context. For the
forward introduction of $L(E)$, the retained $E$ proves its successor
row, independently of the progress assumption and of the chosen witness.
For the reverse direction, open $L(E)$, prove its zero row, and Use at $1$.

For equations $U=(s=t)$ and $V=(r=q)$ define
\[
 U\Rightarrow_g V
 =D_{a,w}L\bigl(C(w,s,r)=C(w,t,q)\bigr).                  \tag{3}
\]
Its zero endpoint is interderivable with $U$, and every successor
endpoint with $V$. The outer witness coordinate is unused. It follows
from Intro, Open, Use, (2), and Lemma~\ref{lem:hcut} that (3) has the
ordinary introduction and elimination rules for an implication between
equations, in hereditary proof families.

Numerical cases are likewise available in hereditary families. For a
compound $A(y)$, introduce $D_{a,y}A(y)$ with unused $a$, ignoring its
base assumption in the supplied successor branch; open it and use the
given base at the desired target. For an equation $E(y)$ apply the
compound construction to $L(E(y))$ and use (2). These constructions do
not require any name theorem or chosen existential witness.

\subsection{Embedding a source equational calculus}

Let $\EE$ be ordinary equational reasoning on $0,S,I$, with the two
computation equations and contextual induction for one literal equation:
\[
 \frac{\Gamma\vdash E(0)\qquad\Gamma,E(z)\vdash E(Sz)}
      {\Gamma\vdash E(t)},
 \qquad z\notin\FV(\Gamma,t).
                                                               \tag{EI}
\]
The eigenparameter $z$ is also absent from the fixed parameters of $E(y)$.
The source calculus has no coverage, separation, or non-equational induction axiom.
Its assumptions are equations. Equality rules include congruence and
substitution in finite term templates.

\begin{lemma}[Equational embedding before names]
\label{lem:embed}
For every finite $\EE$ derivation $\Gamma\vdash E$ and every numerical
term substitution $\sigma$, its instance has a hereditary target proof:
\[
                       \Gamma\sigma\Vdash E\sigma.
\]
\end{lemma}
\begin{proof}
Induct on the source derivation, simultaneously for all substitutions
and receiving contexts. Computation and equality use the hereditary
algebraic rules and Lemma~\ref{lem:hcut}. Here reflexivity itself is
derived: if $r=I(0,t;u.0)$, use the equation $r=t$ as both Transport
premises with template $v=t$ to obtain $t=t$.

Consider (EI), after substitution of its fixed parameters; write its
invariant as $E(y)$ and let $J(y)=G_yE(y)$. By the source induction
hypothesis, the proper step subderivation can be regenerated at
$y=0$ and at $y=Sw$, in arbitrary receiving scopes. Thus, with the
fixed outer context retained, there are hereditary proofs
$E(0)\vdash E(1)$ and $E(Sw)\vdash E(SSw)$. This does not substitute
into a completed target proof.

Put $\Phi(y)=J(y)\Rightarrow_g E(Sy)$. At $y=0$, its implication
proof ignores $J(0)$ and uses the source base and regenerated zero
step. At $y=Sw$, remove the successor guard from $J(Sw)$ and use
the regenerated successor step. Hereditary numerical cases yield
$\Phi(z)$ under the outer context, which is independent of $z$.
Under $J(z)$, implication elimination yields $E(Sz)$ and hence
$J(Sz)$. Guarded progress introduction with unused witness coordinate
now proves $D_{a,y}J(y)$. Open it; its zero row is a computation
theorem. Use at $Sw$ and remove the successor guard to get $E(Sw)$.
Numerical cases with the original base give $E(t)$.

All recursive proof-generation calls in this argument concern proper
subderivations of the source proof. Every target replacement has an
already hereditary donor. This establishes a well-founded construction
and completes the induction.
\end{proof}

\begin{proposition}[Predecessor before names]
\label{prop:pred}
There is a fixed finite raw scalar term $P$ such that, hereditarily
and under all term substitutions,
\[
                 P(0)=0,\qquad P(St)=t.
\]
In particular $Sb=St$ implies $b=t$ by equality reasoning.
\end{proposition}
\begin{proof}
Appendix~\ref{sec:equational-predecessor} constructs $P$ in $\EE$,
using only computation, equality, and equation induction. Apply
Lemma~\ref{lem:embed}. The successor-injectivity conclusion follows
by applying $P$ to the assumed equation and using the two instantiated
predecessor equations.
\end{proof}

The equational construction is important. Using the previously established
whole iPRA equivalence to obtain predecessor here would make the naming
bootstrap circular. Appendix~\ref{sec:equational-predecessor} instead
removes the Boolean case distinctions from the inherited scalar
construction by an explicit equation-induction argument.

\section{Uniform naming with no syntactic exception}
\label{sec:naming}

Define the canonical packet
\[
 A(r):=D_{x,v}\bigl(
 C(v,0,C(r,0,Sx))=C(v,0,r)\bigr).                        \tag{4}
\]
Both sides of its atomic matrix are iterators with counter $v$ and
initializer zero, so (4) satisfies Guard. Its ordinary meaning is
\[
 A(r)^\circ\ \longleftrightarrow\
          \exists x\bigl(C(r,0,Sx)=r\bigr).              \tag{5}
\]
The equivalence follows from its computationally true zero row, its
row at $1$, and numerical coverage in the ordinary interpretation.
Equation (5) explains the packet; the syntactic proof below does not
assume coverage.

\begin{lemma}[Uniform coverage packet]
\label{lem:packet}
For unused $a$, $D_{a,r}A(r)$ has a hereditary $\Gn$ proof. In every
receiving scope and context, $A(St)$ has an effectively generated proof
for every well-scoped term $t$.
\end{lemma}
\begin{proof}
First, $A(0)$ can be introduced at whatever witness (1) selects.
Its successor row computes to $0=0$, independently of that witness
and of its local progress assumption.

To introduce $D_{a,r}A(r)$, its matrix is compound, so Intro asks for
$A(Sz)$ under its base assumption, with a newly allocated newest
parameter $z$. Introduce $A(Sz)$ in that branch. The parent scope of
this inner introduction ends in $z$, so its witness is exactly $z$.
The successor row computes to $Sz=Sz$ and does not use its progress
assumption. The outer base assumption can be ignored. This builds the
whole outer proof in any receiving scope.

The whole construction is hereditary because it allocates $z$ anew on
each regeneration. We do not claim that its isolated internal proof of
$A(Sz)$ can be replayed after $z$ ceases to be newest.

Open $D_{a,r}A(r)$. The dummy witness is absent from the clause
$\langle A(0);r.A(r)\rangle$. Regenerate the first construction of
$A(0)$ in that enlarged scope and Use the clause at $St$. The result
$A(St)$ contains no dummy witness, so Open discharges it.
\end{proof}

\begin{lemma}[Alias extraction]
\label{lem:alias}
For any outer term $t$, open $A(St)$ at a fresh newest parameter $b$.
In that opened context, $b=t$ has a proof that can be regenerated in
every later extension of the numerical scope.
\end{lemma}
\begin{proof}
The opened clause has a zero row that is a computation theorem. Use
it at $1$. Two successor computations give $Sb=St$. Apply
Proposition~\ref{prop:pred}. The retained clause and the hereditary
predecessor instances permit this extraction in any later scope,
including below a new step parameter. It is unnecessary to weaken a
previously completed extraction proof.
\end{proof}

\begin{corollary}[The old literal naming theorem is derived]
\label{cor:name}
For every term $t$, the formula
\[
 D_{x,y}\bigl(C(y,0,x)=C(y,0,t)\bigr)                     \tag{6}
\]
has a hereditary $\Gn$ proof, with no pattern-based witness selection.
\end{corollary}
\begin{proof}
Generate $A(St)$ and open it at newest $b$. Extract $b=t$.
Introduce (6); its selected witness is now $b$. In its step branch,
regenerate the alias equality and use the selector computations to
prove the successor row, ignoring the progress hypothesis. Discharge
the packet opening. Every fresh witness is regenerated inside the
construction, so the resulting proof family is hereditary.
\end{proof}

\section{Completeness and both interpretations}
\label{sec:completeness}

\subsection{Recovering every unrestricted introduction}

We prove the nontrivial implication of Theorem~\ref{thm:main} in a
strong form: every $\G$ derivation gives a hereditary $\Gn$ proof of
the same sequent. Induct on the source derivation, universally over
injective parameter embeddings and receiving contexts. Computations,
Use, and Transport are copied with their parameters. For Open, allocate
a new receiving parameter and apply the induction hypothesis to the
proper branch derivation with that extended embedding.

Consider a source introduction at an arbitrary term $t$, with matrix
$B$ and step derivation
\[
          \Gamma,B(t,\eta_B(z))\vdash_{\G}B(t,Sz).       \tag{7}
\]
In the receiving scope, generate $A(St)$ and open it at a new newest
$b$. Introduce the desired $D_{x,y}B(x,y)$ there. This introduction
selects $b$, and allocates its new step parameter $z$ only afterwards.
In that branch, Lemma~\ref{lem:alias} regenerates $b=t$. Transport
turns the local assumption $B(b,\eta_B(z))$ into
$B(t,\eta_B(z))$. The source induction hypothesis translates the
proper derivation (7) under the enlarged parameter embedding and
context. Lemma~\ref{lem:hcut} supplies its required ordinary assumption
from the hereditary transported assumption. The result is $B(t,Sz)$;
transport back along the alias gives $B(b,Sz)$.

Close the deterministic introduction and the alias opening. The goal
has no $b$, so the eigenparameter condition holds. The two modes agree
because $\eta_B$ depends on the matrix root, which the parameter
embedding and witness substitution do not change. All recursive calls
are to proper source subderivations; the packet and alias proofs were
constructed independently beforehand. This proves $\G\subseteq\Gn$
and finishes Theorem~\ref{thm:main}, using the incoming equivalence
$\Gd=\G$~\cite{deterministic}.

\begin{remark}[Structural rules after completeness]
The theorem now gives ordinary weakening, injective renaming, cut,
and numerical substitution as admissible consequence transformations,
by translation through $\G$. It does not say that an unchanged proof
tree passes the checker after a numerical scope extension.
\end{remark}

\subsection{Ordinary realization and full iPRA}

The ordinary reading is
\[
 (s=t)^\circ=(s=t),\qquad
 (D_{x,y}B)^\circ=
       \exists x\bigl(B^\circ(x,0)\to\forall y\,B^\circ(x,y)\bigr).
                                                               \tag{8}
\]
Clause readings are conditional universal statements as above. The
domain formula is $x=x$, and numerical equality is actual equality.
The guarded zero row makes the atomic introduction in (8) sound by
equation induction; the compound introduction is sound by numerical
zero/successor coverage. Open, Use, and Transport have their usual
intuitionistic readings.

For clarity, the inherited ordinary arithmetic theory $T_I$ is
intuitionistic first-order logic with equality, scalar iterators with
their two equations, contextual induction for literal equations, and
\[
 \forall x\,(x=0\lor\exists y\,x=Sy),\qquad
                  (0=S0)\to\bot.                       \tag{9}
\]
The falsity bridge in (9) is explicit. A native presentation with
designated contradiction $0=S0$ must not silently claim that positive
equations alone exclude the singleton model.

The existing guarded translation from the unrestricted Choice-DI
calculus pads equation matrices. Let $P_\mathrm{fm}(E)=L(E)$ at
equations and $P_\mathrm{fm}(A)=A$ at compound formulas. Its formula
translation is
\[
 T(s=t)=s=t,\qquad T(D_{x,y}B)=D_{x,y}P_\mathrm{fm}(TB).   \tag{10}
\]
Literal inclusion is the reverse translation. The inherited padding
and binder-congruence proofs give both formula round trips uniformly
under assumptions~\cite{guarded}. Theorem~\ref{thm:main} transfers
these very same results to $\Gn$; it is equality of open consequence
relations, not merely agreement on closed arithmetic sentences.

The remaining inherited arithmetic bridge constructs a genuine
predecessor, equality reflection, quantifier-free induction, pairing,
and a compiler for every fixed primitive-recursive description in that
order~\cite{bridge,lean}. In one direction, interpret a scalar iterator
by primitive recursion with an update that ignores the counter. In the
other, if $\pi,L_0,L_1$ are the constructed pairing and projections,
the recursion
\[
 f(0,\vec p)=g(\vec p),\qquad
 f(Sn,\vec p)=h(n,f(n,\vec p),\vec p)
\]
is compiled through
\[
 \begin{aligned}
 V(n,\vec p)&=I\bigl(n,\pi(0,g^*(\vec p));
          u.\pi(SL_0(u),h^*(L_0(u),L_1(u),\vec p))\bigr),\\
 f^*(n,\vec p)&=L_1(V(n,\vec p)).
 \end{aligned}                                         \tag{11}
\]
This is external recursion on a finite function description; there is
no universal total evaluator quantified over all primitive-recursive
codes. Equation induction establishes the compiled defining equations.

For both interpretation composites, internal uniqueness of iteration
and primitive recursion identifies the double-translated operations
with the originals. The same induction over finite formula syntax,
using (8), (10), and the ordinary logical macros, identifies every
formula with its round trip. No quotient or restriction to standard
numerals occurs.

\begin{theorem}[Bi-interpretability and equiconsistency]
The ordinary realization presented by $\Gn$ and full first-order
$\ipra$ are bi-interpretable. In both directions the numerical domain
is whole, equality is actual equality, and the graph of the
comparison isomorphism is $x=y$. Both formula composites are provably
equivalent to the original formulas. Effective translations of finite
proofs preserve contradiction in both directions, so the theories
are equiconsistent.
\end{theorem}
\begin{proof}
Compose Theorem~\ref{thm:main}, the guarded/Choice-DI translations
and their formula round trips, and the native/$T_I$/$\ipra$ bridge.
The first equivalence is identity on syntax and numerical domain;
the inherited interpretations have the properties just described.
Their composite therefore has those properties as well. The explicit
falsity bridge and designated contradiction give the consistency
equivalence. This is an external metatheorem, with no claim that
either theory proves its own consistency.
\end{proof}

\section{Executable presentation and verification}
\label{sec:implementation}

\subsection{Certificate grammar}

The executable checker retains the incoming strict format. The JSON
object has exactly the keys \code{goal} and \code{proof}. Numerical
syntax uses integer zero, positive de Bruijn indices for bound
variables, and negative levels for free parameters. Successor is
\code{[0,t]}, iteration is \code{[n,a,h]}, an equation is \code{[s,t]},
and a paired formula is \code{[B]}. The scopes and the sorts distinguish
the two uses of binary lists. The iterator body binds one variable;
the paired matrix binds two.

The full instruction inventory is
\begin{center}\small
\begin{tabular}{ll}
\toprule
Instruction & Contents\\\midrule
\code{zero} & initializer, update\\
\code{succ} & counter, initializer, update\\
\code{use} & clause index, target term, antecedent proof\\
\code{tr} & formula template, equation proof, instance proof\\
\code{open} & paired major proof, branch proof\\
\code{intro} & matrix, branch proof\\
\bottomrule
\end{tabular}
\end{center}
Every constructed output is checked for sort, guard, and scope. Inputs
use exact integer types, so Boolean values are rejected as numerical
syntax. A branch cannot export its local parameter even in an unused
iterator update. Python's optimization flag does not disable any
acceptance-critical check; the program does not use \code{assert} for
input validation.

The sole change to the incoming checker is
\begin{lstlisting}
-        t = witness(matrix, scope)
+        t = -scope
\end{lstlisting}
together with deletion of the now unused \code{witness} function.
At scope zero, \code{-scope} is literal zero; otherwise it names the
newest parent parameter. The fresh step parameter remains
\code{-scope - 1}. No part of the new proof construction is hidden
inside the verifier.

\subsection{Measurements and limits}

The measurement script removes actual Python docstrings, applies
\code{ast.unparse}, and adds one terminal newline. Bytes are UTF-8
bytes of this normalized complete module. AST size counts all nodes
from \code{ast.walk}. Significant tokens exclude encoding, end marker,
indentation, line separators, and comments by their named
\code{tokenize} constants. The results reproduce
Table~\ref{tab:main-size} on CPython 3.12.14.

The reductions are $6.75\%$ in bytes, $7.42\%$ in tokens, and
$6.80\%$ in AST nodes. The comparison programs have identical parser,
public language, error handling, computation rules, transport rule,
and scope discipline. This is an axiomatization change with a measured
implementation benefit, rather than a serialization-only change.

The supplied validation scripts run 11 accepted and 17 rejected certificates
in both normal and optimized Python: 56 CLI checks, all passing. They
exercise accepted derivations, malformed
syntax, guards and binders, eigenparameter escape, the whole universal
coverage construction, and the exact syntactic patch against the frozen source control. A
separate comparison confirms that the two witness policies need different
branch certificates even though their theorem sets coincide. All eleven
positive certificates are also elaborated by inserting their selected
witnesses and replayed successfully in the frozen unrestricted source
checker. Exact receipts are
generated by \code{reproduce.py}; they distinguish executed finite
checks from the universal written proofs in this paper.

The implementation does not include a complete executable compiler
from every unrestricted proof to the new calculus. That transformation
is proved in Sections~\ref{sec:hereditary}--\ref{sec:completeness}.
In particular, the full predecessor derivation and the general
hereditary-proof metatheorem have not been checked by Lean or Rocq in
this package. Existing upstream Lean theorems are credited foundations,
not a claim that the newly restricted calculus is already formalized.
Large finite certificates can also exceed a particular Python process's
stack or memory limits; no uniform resource bound is asserted.

\section{What this settles and what remains open}

The current raw deterministic presentation does not require a special
case for literal naming matrices. Arbitrary names can be generated
using its ordinary numerical parameter supply. This removes a branch
from the mathematical specification and 249 normalized bytes from its
complete strict checker.

The result also clarifies two dependencies. The inherited scalar
predecessor construction can be carried out in the purely equational
fragment; numerical coverage is not necessary for its two defining
equations. Conversely, a proof that uses the newest parameter is not
automatically stable when another parameter is appended. Hereditary
regeneration is the reason that the naming and completeness arguments
remain valid.

An absolute smallest presentation is not established. It would require
a fixed comparison class for languages, syntax checking, certificate
protocols, libraries, and source-size conventions, together with a
lower-bound argument. The compact codecs and the newer three-family
calculi remain relevant alternative implementations. The present
contribution is the precise deletion and equivalence theorem above,
with its complete, reproducible checker comparison.

\appendix
% Research note for the parent agent. No external logic or coverage is used.
\section{Predecessor without numerical coverage}
\label{sec:equational-predecessor}

Let $\mathcal E$ be the purely equational calculus on $0,S,I$, with
ordinary equality rules, the two iterator equations, and literal-equation
induction with parameters.  In particular, $\mathcal E$ has no axiom saying
that every number is zero or a successor.  The following strengthens the
dependency statement of the inherited scalar predecessor construction:
\[
  \mathcal E\vdash P(0)=0,
  \qquad \mathcal E\vdash P(Sx)=x
\]
for one finite iterator term $P$.  Thus $\mathcal E$ proves successor
injectivity.  The construction below is inherited; the new point is the
explicit elimination of all its numerical and Boolean case splits.

Write $1=S0$ and define
\[
 C(e,a,b)=I(e,a;u.b),\qquad Z(e)=C(e,1,0),\qquad
 a+b=I(b,a;u.Su).
\]
The equations of a commutative monoid for addition, and
$Sa+b=S(a+b)$, follow by equation induction.  Expressions such as $2a$ and
$3a$ abbreviate repeated addition.

\paragraph{Equational interpolation.}
For every term template $F(v)$, with all other variables fixed, equation
induction on a fresh $e$ proves
\begin{equation}
 F(Z(e))=C(e,F(1),F(0)). \tag{BI}
\end{equation}
At $e=0$ both sides are $F(1)$; at $e=Sz$ both sides are $F(0)$.
The induction hypothesis is unused.  Also equation induction, again with
an unused hypothesis, proves $Z(Z(Z(e)))=Z(e)$.
Consequently, if $b=Z(Z(b))$, and the two equations
$F(0)=G(0)$ and $F(1)=G(1)$ are available under the current fixed
parameters, then $F(b)=G(b)$ follows by applying (BI) at $e=Z(b)$.
This is a derived equality rule, not a disjunction elimination rule.  It
continues to hold in the presence of arbitrary fixed side equations.

\paragraph{Parity and a three-cycle.}
Put
\[
 N(n)=I(n,0;u.Z(u)),\quad j(u)=Z(u)+2N(u),\quad
 M(n)=I(n,0;u.j(u)).
\]
Equation induction proves $N(n)=Z(Z(N(n)))$: the zero case computes;
the successor case is $Z(N(n))=Z(Z(Z(N(n))))$.  Applying (BI) therefore
gives
\[
 N(n)+Z(N(n))=1,\qquad N(n+2)=N(n).
\]
Equation inductions give
\[
 N(2n)=0,\qquad N(2n+1)=1,\qquad N(3n)=N(n).
\]
For the last equation use $N(a+3)=Z(N(a))$ in the successor step.
Finite computations yield $j(0)=1,j(1)=2,j(2)=0$.

\paragraph{The ramp invariant, with its case split removed.}
Define
\[
 G(n)=I(n,0;u.u+M(u)+1).
\]
We prove $M(G(n))=N(n)$ by equation induction.  At the successor step
write $a=G(n)$ and $b=N(n)$, so the sole induction equation is $M(a)=b$.
Interpolation in the Boolean parameter $b$ gives
\[
 M(a+b+1)=C(b,M(a+1),M(a+2)).
\]
Both endpoint equations follow from the computations of $C$.  Rewriting
by $M(a)=b$ and the two computations of $M$ turns the right side into
$C(b,j(b),j(j(b)))$.  A further use of (BI) gives
\[
 C(b,j(b),j(j(b)))=Z(b):
\]
at $b=0$ both sides are $1$, and at $b=1$ both sides are $0$.
Since $G(Sn)=a+b+1$ and $N(Sn)=Z(b)$, this proves the induction step.
No disjunctive fact about $b$ was used.

It follows by equality reasoning and equation induction that
\[
 G(Sn)=G(n)+N(n)+1,\quad G(n+2)=G(n)+3,\quad G(2n)=3n.
\]
Another equation induction proves
\[
 G(3n)=3G(n)+N(n).
\]
Its successor step reduces, with $b=N(n)$, to the Boolean identity
$2b+4=3b+3+Z(b)$, whose endpoints are checked directly and whose
general Boolean instance follows by (BI).

\paragraph{The factor map.}
Set $H(u)=C(N(u),G(u),u)$.  We obtain
\[
 H(2u)=3u,\qquad H(1)=1,\qquad H(3u)=3H(u).
\]
Only the last identity requires comment.  With $b=N(u)$ it is
\[
 C(b,3G(u)+b,3u)=3C(b,G(u),u),
\]
which is again (BI): its zero endpoint is $3G(u)=3G(u)$ and its one
endpoint is $3u=3u$.  The apparent dependence between $b$ and $u$
causes no problem: prove the displayed identity with an independent
Boolean variable, then substitute $N(u)$.

\paragraph{Comparator.}
For any fixed update $h$, equation induction gives rotation
$I(Sy,a;h)=I(y,h(a);h)$.  Equation induction also gives
\[
 I(y,3a;H)=3I(y,a;H),\qquad I(y,1;H)=1.
\]
Let
\[
 E(x)=I(x,1;u.2u),\qquad L(x,y)=N(I(y,E(x);H)).
\]
The preceding identities give the complete comparator specification
\[
 L(0,y)=1,\qquad L(Sx,0)=0,\qquad L(Sx,Sy)=L(x,y).
\]
For the third equation the calculation is
\[
\begin{aligned}
 N(I(Sy,2E(x);H))
 &=N(I(y,H(2E(x));H))\\
 &=N(I(y,3E(x);H))\\
 &=N(3I(y,E(x);H))=N(I(y,E(x);H)).
\end{aligned}
\]

\paragraph{Saturation and predecessor.}
Put
\[
 B_m(u)=C(L(Su,m),u,Su),\qquad W(t,m)=I(t,0;u.B_m(u)).
\]
The identity $C(e,Sa,Sb)=SC(e,a,b)$ is an equation induction on $e$
with unused induction hypothesis.  It gives
\[
 B_{Sm}(Su)=SB_m(u),\qquad B_{Sm}(0)=S0.
\]
Equation induction on $t$, with $m$ fixed, gives
$W(St,Sm)=SW(t,m)$.  Equation inductions then give
\[
 L(Sn,n)=0,\qquad W(n,n)=n,\qquad W(Sn,n)=n.
\]
In the middle equation substitute $t=m=n$ into the already established
two-parameter shift equation.  No parameter of a live induction hypothesis
is changed.

Finally define the finite scalar term
\[
 P(m)=I(m,0;u.C(L(SSu,m),u,Su)).
\]
Its zero value computes.  The comparator recurrence identifies the update
of $P(Sm)$ with $B_m(u)$.  Formally, equation induction on a fresh counter
$t$ proves
\[
 I(t,0;u.C(L(SSu,Sm),u,Su))=W(t,m).
\]
At $t=Sm$ the right side is $m$, proving $P(Sm)=m$.  This last
induction accounts for iterator-body extensionality; it is not an extra
equality axiom.  Applying $P$ to $Sx=Sy$ proves $x=y$.

\paragraph{Audit boundary.}
The proof uses only the two iterator equations, equality reasoning,
substitution into established equations, and equation induction with fixed
side parameters.  Every Boolean step has been replaced by (BI).  It does
not use coverage, a predecessor axiom, successor injectivity, equality
reflection, pairing, or induction on a non-equational formula.  This is a
written syntactic derivation; it is not a new Lean verification.


\section{Reproduction and provenance}

From this directory run
\begin{lstlisting}
python3 reproduce.py
python3 code/checker.py < code/examples/accepted/universal_coverage_closure.json
pdflatex -no-shell-escape -halt-on-error paper.tex
pdflatex -no-shell-escape -halt-on-error paper.tex
\end{lstlisting}
The exact example inventory and test commands are recorded in
\code{README.txt}. The package includes the frozen incoming checker,
the new checker, normalized sources, measurements, validation scripts,
the mathematical notes, and the inherited bridge sources used in the
proof. \code{PROVENANCE.json} identifies the repository snapshot and
the five incoming archives; \code{SHA256SUMS} covers the delivered
payload. The upstream material is distributed under its MIT No
Attribution license, preserved in \code{LICENSE-upstream.txt}.

\begin{thebibliography}{9}
\raggedright
\bibitem{repo} ProveIt contributors.
\emph{Concise presentations of intuitionistic primitive recursive arithmetic}.
Snapshot \code{9f9545f5ffd9e0392d4a36639503aca12383450a},
November 12, 2026.
\url{https://github.com/VladimirReshetnikov/ProveIt/tree/9f9545f5ffd9e0392d4a36639503aca12383450a/Logic/PrimitiveRecursiveArithmetic/ConciseIPRA}.
\bibitem{deterministic} \emph{Deterministic introduction for intuitionistic
primitive recursive arithmetic}. Incoming archive
\code{ipra\_deterministic\_introduction.zip}, in the snapshot above.
Its strict checker is the matched control in this package.
\bibitem{parameter} \emph{iPRA without supplied introduction witnesses}.
Incoming archive \code{ipra\_parameter\_intro\_2026-11-11.zip},
in the snapshot above.
\bibitem{evaluation} \emph{iPRA with evaluation at witness discharge}.
Incoming archive \code{ipra\_evaluation\_discharge\_2026-11-11.zip},
in the snapshot above.
\bibitem{guarded} ProveIt, \code{Research/guarded\_core} and its
\code{guarded\_core.tex} bridge, also reproduced in the incoming
deterministic-introduction archive. Its canonical formula grammar,
padding, equation implication, and guarded/native translations are
inherited here.
\bibitem{bridge} ProveIt, \code{Research/certified\_zero/reference}:
\code{inherited\_choice\_basis.tex},
\code{inherited\_ordinary\_bridge.tex},
\code{scalar\_predecessor.tex},
\code{equality\_reflection.tex}, and
\code{primitive\_recursion.tex}. Copies are included under
\code{reference/}.
\bibitem{lean} ProveIt, \code{Lean/ConciseIPRA}, especially
\code{NativePublic.lean}, \code{NativePublicD.lean},
\code{NativeBiInterpretation.lean}, and \code{BiInterpretation.lean}.
The snapshot records the established native/scalar/iPRA equivalences;
it does not formalize the new restriction proved here.
\end{thebibliography}
\end{document}
